
Errors that are most difficult to repair often provide the least actionable feedback. Assertion errors constitute 53.4\% of failures on standard code generation benchmarks, yet their error messages contain near-zero extractable state information. This creates a fundamental tension in LLM-based code repair: the errors most in need of rich diagnostic feedback are those that provide the least.

Recent advances in self-repair demonstrate that iterative refinement with verification feedback can improve pass rates by 4.9--17.1 percentage points on HumanEval. Runtime traces outperform simple error messages, and type constraints improve functional correctness. However, these findings describe which feedback works better without explaining why. The question remains: what properties of a verification signal enable an LLM to localize bugs and infer correct fixes?

This paper hypothesizes that Actionable Specificity (AS) mediates feedback effectiveness. AS is decomposed into three measurable components: localization (AS\_loc), indicating whether file:line information is present; state exposure (AS\_state), counting how many variable values are visible; and causal context (AS\_causal), measuring execution trace depth between assignment and failure. Under this framework, a full runtime trace (high AS) should outperform a bare assertion error (low AS) because it provides localization, state, and causal information necessary for targeted repair.

To test whether AS is measurable from standard verification signals---a necessary precondition for testing its predictive power---an existence test was designed on HumanEval and MBPP (664 problems). 600 signals were generated across 6 conditions (C1--C6) varying in expected AS level, and AS components were extracted using regex patterns designed for Python tracebacks.

The key finding is negative: regex-based extraction achieves only 49.5\% coverage, with AS\_state at 2.7\%. The root cause is that assertion errors, which dominate these benchmarks (53.4\%), produce minimal output without traceback frames. Standard AssertionError messages lack the file:line and variable context that extraction patterns require.

This limitation of the operationalization is acknowledged. However, the failure is informative: it reveals a boundary condition that future work must address. The contributions are:

\begin{enumerate}
\item The AS framework provides a principled decomposition for analyzing feedback informativeness.
\item Error type stratification is identified as a necessary precondition: feedback informativeness research must verify extraction feasibility by error type before operationalizing.
\item AS ordering ($C1\geq C2\geq C3\geq C4)$ is preserved on the extractable subset, suggesting the conceptual decomposition may remain valid under alternative operationalizations.
\item The framework provides actionable guidance by identifying which error types block extraction and why.
\end{enumerate}

